\subsection{Данные экспериментального исследования}

\subsubsection{Описание сценариев нагрузки}

Для сравнительного исследования производительности двух реализаций
пула потоков были разработаны два сценария нагрузки, отражающих
типичные паттерны использования пулов потоков в производственных
системах. Каждый сценарий реализован в виде параметризованного
бенчмарка библиотеки Google Benchmark, запускаемого с числом
рабочих потоков 1, 2, 4, 8 и 16.

Первый сценарий --- много мелких задач --- моделирует ситуацию,
характерную для высоконагруженных серверных систем, в которых
пул потоков обрабатывает большое число коротких независимых
запросов. В данном сценарии в пул последовательно ставятся
10\,000 задач, каждая из которых выполняет атомарный инкремент
счётчика. Атомарный инкремент выбран намеренно: его время
выполнения составляет единицы наносекунд, что обеспечивает
доминирование накладных расходов планировщика над полезной
работой. Таким образом, данный сценарий измеряет прежде всего
эффективность механизма постановки и извлечения задач, а не
производительность самих задач. Реализация сценария представлена
в листинге~\ref{lst:bench-small}.

\begin{lstlisting}[language=C++,
                   caption={Сценарий 1: много мелких задач},
                   label={lst:bench-small}]
static void BM_ManySmallTasks(benchmark::State& state,
                               auto& pool)
{
    for (auto _ : state) {
        std::atomic<int> counter(0);
        std::vector<std::future<void>> futures;
        futures.reserve(kTaskCount); // kTaskCount = 10000

        for (int i = 0; i < kTaskCount; ++i) {
            futures.emplace_back(pool.submit([&counter]() {
                counter.fetch_add(1,
                    std::memory_order_relaxed);
            }));
        }

        for (auto& f : futures) { f.get(); }
    }

    state.SetItemsProcessed(
        state.iterations() * kTaskCount);
}
\end{lstlisting}

Для данного сценария ожидается следующее поведение. При малом
числе потоков обе реализации должны демонстрировать схожую
производительность, поскольку конкуренция за разделяемые
структуры данных невелика. При увеличении числа потоков
\texttt{SimpleThreadPool} должен деградировать быстрее из-за
конкуренции всех потоков за единственную глобальную очередь.
\texttt{WorkStealingPool} теоретически должен масштабироваться
лучше, однако использование мьютекса в \texttt{WorkStealingDeque}
ограничивает это преимущество.

Второй сценарий --- неравномерная нагрузка --- моделирует
ситуацию, в которой задачи существенно различаются по времени
выполнения. Данный паттерн характерен для систем обработки
запросов, где часть запросов требует обращения к базе данных
или внешним сервисам, а другая часть обрабатывается локально.
В данном сценарии в пул ставятся 10\,000 задач, половина из
которых является короткими (атомарный инкремент), а другая
половина --- длинными (задержка 200~мкс посредством
\texttt{std::this\_thread::sleep\_for}). Короткие и длинные
задачи чередуются: чётные по индексу задачи являются длинными,
нечётные --- короткими. Реализация сценария представлена
в листинге~\ref{lst:bench-uneven}.

\begin{lstlisting}[language=C++,
                   caption={Сценарий 2: неравномерная нагрузка},
                   label={lst:bench-uneven}]
static void BM_UnevenLoad(benchmark::State& state,
                           auto& pool)
{
    for (auto _ : state) {
        std::atomic<int> counter(0);
        std::vector<std::future<void>> futures;
        futures.reserve(kTaskCount);

        for (int i = 0; i < kTaskCount; ++i) {
            futures.emplace_back(
                pool.submit([&counter, i]() {
                    if (i % 2 == 0) {
                        // Длинная задача: 200 мкс
                        std::this_thread::sleep_for(
                            std::chrono::microseconds(200));
                    }
                    counter.fetch_add(1,
                        std::memory_order_relaxed);
                }));
        }

        for (auto& f : futures) { f.get(); }
    }

    state.SetItemsProcessed(
        state.iterations() * kTaskCount);
}
\end{lstlisting}

Для данного сценария ожидается что \texttt{WorkStealingPool}
продемонстрирует преимущество перед \texttt{SimpleThreadPool}
при большом числе потоков. Механизм кражи задач позволяет
незагруженным потокам забирать задачи у потоков, занятых
выполнением длинных задач, что должно обеспечить более
равномерную загрузку всех ядер процессора.

Выбор данных двух сценариев обусловлен тем, что они охватывают
два полярных случая с точки зрения балансировки нагрузки.
В первом сценарии все задачи одинаковы по длительности ---
балансировка нагрузки не требуется и оба пула должны работать
схожим образом. Во втором сценарии задачи существенно различаются
по длительности --- именно здесь механизм work-stealing имеет
теоретическое преимущество. Сравнение результатов двух сценариев
позволяет выявить условия, при которых механизм work-stealing
оправдывает свои накладные расходы.

Параметры бенчмарков представлены в таблице~\ref{tab:bench-params}.

\begin{table}[h]
\caption{Параметры экспериментального исследования}
\label{tab:bench-params}
\begin{tabularx}{\textwidth}{|l|X|}
\hline
\textbf{Параметр} & \textbf{Значение} \\
\hline
Число задач в сценарии & 10\,000 \\
\hline
Длительность короткой задачи &
Атомарный инкремент ($\approx$ 5~нс) \\
\hline
Длительность длинной задачи &
\texttt{sleep\_for(200~мкс)} \\
\hline
Исследуемые значения числа потоков &
1, 2, 4, 8, 16 \\
\hline
Число итераций & Определяется автоматически
                 библиотекой Google Benchmark \\
\hline
Режим измерения времени & Реальное время (\texttt{UseRealTime}) \\
\hline
Уровень оптимизации компилятора & \texttt{-O2} \\
\hline
\end{tabularx}
\end{table}

\subsubsection{Условия проведения эксперимента}

Корректность сравнительного исследования производительности
требует фиксации условий проведения эксперимента. В данном
разделе описываются характеристики тестовой машины и параметры
запуска бенчмарков.

Эксперимент проводился на рабочей станции под управлением
операционной системы Linux. Характеристики аппаратного
обеспечения были получены непосредственно из вывода библиотеки
Google Benchmark и представлены в таблице~\ref{tab:hardware}.

\begin{table}[h]
\caption{Характеристики тестовой машины}
\label{tab:hardware}
\begin{tabularx}{\textwidth}{|l|X|}
\hline
\textbf{Параметр} & \textbf{Значение} \\
\hline
Число логических ядер & 16 \\
\hline
Тактовая частота & 5102,71 МГц \\
\hline
Кэш L1 (данные) & 32 КиБ $\times$ 8 \\
\hline
Кэш L2 & 1024 КиБ $\times$ 8 \\
\hline
Кэш L3 & 16384 КиБ $\times$ 1 \\
\hline
Операционная система & Linux \\
\hline
Компилятор & Clang, стандарт C++17 \\
\hline
Уровень оптимизации & \texttt{-O2} \\
\hline
\end{tabularx}
\end{table}

Следует отметить что при проведении эксперимента была включена
динамическая регулировка тактовой частоты процессора (CPU
frequency scaling). Google Benchmark выводит предупреждение
о данном факте:

\begin{lstlisting}[caption={Предупреждение Google Benchmark
                            о CPU scaling},
                   label={lst:cpu-scaling}]
***WARNING*** CPU scaling is enabled, the benchmark real time
measurements may be noisy and will incur extra overhead.
\end{lstlisting}

Динамическая регулировка тактовой частоты означает что процессор
может изменять частоту в ходе эксперимента в зависимости от
тепловой нагрузки и энергопотребления. Это вносит дополнительный
разброс в результаты измерений. Отключение данной функции
требует прав администратора и изменения системных параметров
ядра Linux, что не всегда возможно в стандартной рабочей среде.
Для компенсации данного эффекта Google Benchmark автоматически
увеличивает число итераций до достижения статистически
значимых результатов.

Параметры запуска бенчмарков определялись следующим образом.
Число итераций для каждого бенчмарка определялось автоматически
библиотекой Google Benchmark: библиотека увеличивает число
итераций до тех пор пока суммарное время выполнения не превысит
пороговое значение, обеспечивающее статистически значимый
результат. Для каждого бенчмарка измерялось реальное время
выполнения (\texttt{real\_time}), а не процессорное время
(\texttt{cpu\_time}), поскольку реальное время отражает
фактическую пропускную способность системы с учётом времени
ожидания потоков.

Исследуемые значения числа потоков --- 1, 2, 4, 8 и 16 ---
выбраны как степени двойки в диапазоне от одного потока до
полного использования всех логических ядер тестовой машины.
Данный диапазон позволяет исследовать три характерных режима
работы пула потоков. При числе потоков меньше числа ядер
(1, 2, 4) система недозагружена и накладные расходы планировщика
минимальны. При числе потоков равном числу ядер (16) достигается
теоретически оптимальная загрузка. Промежуточное значение 8
позволяет наблюдать переходный режим.

Каждый бенчмарк запускался в рамках единственного процесса
без параллельного выполнения других бенчмарков, что исключает
взаимное влияние сценариев нагрузки на результаты измерений.
Бенчмарки и тесты корректности запускались раздельно: тесты
собираются с Thread Sanitizer, который вносит значительные
накладные расходы и не должен влиять на измерения
производительности.